\subsection{Two eight-element triangular extensions with a low-degree quotient}

\begin{theorem}[Proved: the eight-element sources seven and eighteen]
\label{CP2:A-eight-triangular-low-degree-quotients}
Let $P=\{1,e_{ij}:i,j\in\{0,1\}\}$ be the
identity-adjoined rectangular band and write $e=e_{00}$.
Adjoin three elements $a,b,c$, each with the constant
band actions
\[
 te_{ij}=e_{0j},\qquad e_{ij}t=e_{i0}
       \qquad(t=a,b,c).
\]
Consider either of the two outside multiplication tables
\[
\begin{array}{c|ccc}
 &a&b&c\\\hline
a&e&e&e\\b&e&e&e\\c&a&b&c
\end{array}
\qquad\text{or}\qquad
\begin{array}{c|ccc}
 &a&b&c\\\hline
a&e&e&e\\b&e&e&b\\c&a&b&c.
\end{array}                                                     \tag{1}
\]
For every finite left-$N$, right-$M$ biset $S$, the
monoid retraction $\rho:M\to P$ sending $a,b,c$ to
$e$ extends to an actual quotient functor of collages
\[
 \mathcal C_S\longrightarrow\mathcal C_U
\]
that induces an isomorphism in rational homology and
cohomology through degree two.  It leaves $N$ unchanged
and replaces $S$ by an explicitly defined quotient $P$-set.
If $N=R_3$, all rational degree-two products of
$B\mathcal C_S$ vanish, for every state count and
without a boundedness hypothesis.
The two tables in (1) are respectively positive
order-eight catalog sources $7$ and $18$.
\end{theorem}
\begin{proof}
Put $R=\mathbb QM$, $B=\mathbb QP$ and let
$\Delta$ be the one-dimensional right-$B$ module on
which $1$ acts as one and every band element acts as zero.
Set $h=c-e$ and $p=1-h$.
Literal multiplication in (1) gives
\[
 h^2=h,\quad pRh=0,\quad pRp\cong B,
 \quad hK=Kh=0.
                                                               \tag{2}
\]
Here $\rho$ identifies the corner $pRp$ with $B$,
and its identity $p$ maps to $1$.

In the first case, put $z_a=a-e$, $z_b=b-e$.
Then $Rh=\mathbb Qh$ and the following is an exact
sequence of left-$R$, right-$B$ bimodules:
\[
 0\longrightarrow(Rh\otimes_{\mathbb Q}\Delta)^2
   \xrightarrow{(w_1,w_2)\mapsto w_1z_a+w_2z_b}
 Rp\longrightarrow B\longrightarrow0.                         \tag{3}
\]
In the second case, put $x=b-e$ and $z=a-e$.
Then $Rh=\mathbb Qh\oplus\mathbb Qx$,
$hRh=\mathbb Q[h,x]/(x^2)$ with identity $h$, and
\[
 xz=0,\qquad hz=z,\qquad zK=0,\qquad zp=z.
\]
There is the complete left-projective bimodule resolution
\[
 \cdots\xrightarrow{\,x\,}Rh\otimes\Delta
 \xrightarrow{\,x\,}Rh\otimes\Delta
 \xrightarrow{\,z\,}Rp\longrightarrow B\longrightarrow0.
                                                               \tag{4}
\]
Each displayed arrow is right multiplication by its
label.  Its exactness is literal: multiplication by
$z$ sends $h$ to $z$ and $x$ to zero; multiplication
by $x$ sends $h$ to $x$ and $x$ to zero.
Thus every successive kernel is the preceding image.
The terms $Rh,Rp$ are left projective, and the
external right-$B$ action on every $Rh$ term is
the one-dimensional action $\Delta$.
The identities above make all arrows right-$B$ linear.
The quotient maps in (3) and (4) are the actual
algebra map induced by $\rho$.

Let $L=\mathbb Q[S]$, with its commuting left
$\mathbb QN$ action.  Denote the right actors of
$a,b,c,e$ by $A,D,C,E$, respectively.
The idempotents $C,E$ satisfy $CE=EC=E$.
First form the actual set quotient $T$ by identifying
$C(s)$ with $E(s)$ for every $s$.
It is a left-$N$, right-$P$ set: the two endpoints
have equal images under every band actor, by (1),
and the left action commutes with both idempotents.
As a left-$\mathbb QN$, right-$B$ module,
\[
 Lp\cong\mathbb Q[T].                                          \tag{5}
\]
An explicit integral basis establishing (5) is
\[
 s-C(s)+E(s)\quad(s\notin CS),\qquad t\quad(t\in ES).
\]
It maps bijectively to the state basis of $T$.
Also $Lh$ has basis $t-E(t)$ for $t\in CS\setminus ES$.

In the first case the two maps from $Lh$ to $Lp$
in the tensor of (3) become the graph-incidence vectors
\[
 [A(t)]-[E(t)],\qquad [D(t)]-[E(t)]
       \quad(t\in CS\setminus ES).                            \tag{6}
\]
In the second case the degree-one map in the tensor
of (4) gives just the first vector in (6).
Square brackets mean the state images in $T$.
These identities follow by expanding
$(t-te)(a-e)=ta-te$ and the corresponding expression
with $b$.
Let $U$ be the component set of the finite graph on
$T$ with the edges (6) that apply to the chosen case.
Loops and repeated edges are allowed.
Its linear incidence cokernel is exactly $\mathbb Q[U]$.

This is the required actual quotient set.
The relation is left-$N$ stable by commutation.
It is right-$P$ stable because the two endpoints of
each edge have equal images under every band actor.
All three added actors descend to the action of $e$:
$ca=a$ and $cb=b$ show that their images are already
represented by the states indexed in (6) in the first
case; in the second case $bc=b$ places every $b$-image
in $CS$, where it was collapsed in the first quotient.
The equality $be=e$ identifies its image there with
$E(s)$.  The actor $c$ was likewise collapsed to $e$.
Thus $S\to U$ is equivariant for the literal monoid
retraction $\rho$, and gives an actual collage functor.

Here is the low-degree comparison, including the
unbounded resolution in the second case.
Tensor (3) or (4) on the left with $L$.
The resulting nonnegative right-$B$ complex $F$
computes $L\otimes_R^{\mathbf L}B$.
Its degree-zero term is $Lp$; all positive-degree
terms are copies of $Lh\otimes\Delta$.
The actual quotient induces the natural augmented
chain map
\[
 F\longrightarrow\mathbb Q[U].                                \tag{7}
\]
It is the adjoint of the state-module map
$L\to\mathbb Q[U]$, because its degree-zero map
sends $lp$ to the image of $l$, using $\rho(p)=1$.
The graph-incidence description proves that (7)
is an isomorphism on degree-zero homology.

The kernel complex of (7) has degree-zero term
the incidence image in $Lp$, and has the same
positive-degree terms as $F$.
Every term has right-$B$ action $\Delta$: for the
degree-zero term this follows from $z_aK=z_bK=0$
and $z_ap=z_a$, $z_bp=z_b$ (or the corresponding
single $z$ identities).
Its degree-zero homology is zero.
Truncating this complex by replacing degree one
with the kernel of its degree-one differential
therefore gives a natural quasi-isomorphic complex
of the form
\[
 G\otimes_{\mathbb Q}\Delta[1],                                \tag{8}
\]
where $G$ is a nonnegative complex of left-$\mathbb QN$
modules.  No equivariant splitting and no vanishing
of the higher homology of $G$ is asserted.

The rectangular-band resolution gives
$\Delta\otimes_B^{\mathbf L}\mathbb Q\simeq\mathbb Q[1]$.
Indeed its terms after tensoring are $0,\Delta,0$
in projective degrees zero, one and two.
Consequently the fiber of the actual bridge comparison
after applying the two derived bars is
\[
 \left(\mathbb Q\otimes_{\mathbb QN}^{\mathbf L}G\right)[2].
                                                               \tag{9}
\]
It has zero homology below degree two, since $G$
and the ordinary monoid bar are nonnegative.
This assertion does not require the periodic terms
in (4) to vanish in any degree.

On source corners, the trivial right-$R$ module
annihilates $h$, so tensoring (3) or (4) with it
leaves only $\mathbb Qp=\mathbb Q$ in degree zero.
Derived associativity therefore shows that the actual
map $BM\to BP$ is a rational homology equivalence
in every degree.  The opposite corner is unchanged.
The natural two-block pushout comparison now says
that the mapping cone of the full collage map has
the bridge fiber (9) shifted twice.
It has no homology below degree four.
Thus the actual collage quotient induces an
isomorphism on homology, and dually cohomology,
in degrees zero, one and two, as asserted.

Finally take $N=R_3$.
The trivial left-$B$ resolution has length two and
the trivial right-$\mathbb QN$ resolution has length one.
Hence the quotient collage has rational homology
bounded in degree four, for every finite $U$.
Corollary~\ref{CP2:A-baseline-right-zero-three-six}
therefore kills all its rational degree-two products.
Every second-degree class on the original collage
is the pullback of one on this quotient by the
actual cohomology isomorphism.  Naturality of the
cup product proves the desired vanishing.
\end{proof}
